\documentclass[11pt]{article}

\usepackage[margin=1in]{geometry}
\usepackage{amsmath,amssymb,amsthm,mathtools}
\usepackage{hyperref}
\usepackage{microtype}

\hypersetup{
  colorlinks=true,
  linkcolor=blue,
  citecolor=blue,
  urlcolor=blue,
}

% --- Theorem environments ---
\theoremstyle{plain}
\newtheorem{theorem}{Theorem}
\newtheorem{lemma}[theorem]{Lemma}
\newtheorem{proposition}[theorem]{Proposition}
\newtheorem{corollary}[theorem]{Corollary}
\theoremstyle{definition}
\newtheorem{definition}[theorem]{Definition}
\theoremstyle{remark}
\newtheorem{remark}[theorem]{Remark}

% --- Notation helpers ---
\newcommand{\asinh}{\operatorname{arcsinh}}
\newcommand{\Ci}{\operatorname{Ci}}
\newcommand{\E}{\mathcal{E}}
\newcommand{\dd}{\,\mathrm{d}}

\title{EVQ-Cosh: Theory-Only Notes (Theorems and Key Formulas)}
\author{paperdraft extract from \texttt{paper_draft/CORE\_THEORY.md}}
\date{2026-03-05}

\begin{document}
\maketitle

\begin{abstract}
This note collects the minimal theory required for a paper-ready EVQ-Cosh
presentation: the exact broadband Euler--Lagrange ODE, its closed-form solution,
the EVQ quantile warp (inverse CDF) in closed form, the $\tau\to 0$ geometric
limit, the operator interpretation of the $\min$ kernel, and the Jensen
``waterbed'' inequality. Empirical claims are included only insofar as they
justify the single approximation step (broadband projection).
\end{abstract}

\section{Setup and Notation}

\begin{definition}[RoPE frequency map]
Fix a base $b>1$. We parameterize frequencies by $\phi\in[0,1]$ via
\begin{equation}
\omega(\phi)=b^{-\phi}.
\end{equation}
Thus $\phi=0$ corresponds to the highest frequency ($\omega=1$) and $\phi=1$ to
the lowest frequency ($\omega=b^{-1}$).
\end{definition}

\begin{definition}[Channel density and discretization]
Let $\rho(\phi)\ge 0$ be a density on $[0,1]$ with $\int_0^1\rho(\phi)\dd\phi=1$.
Given $N$ discrete frequency channels, we use the standard RoPE indexing
\begin{equation}
u_k=\frac{k}{N},\qquad k=0,1,\dots,N-1,
\end{equation}
and define discrete locations $\phi_k$ by the inverse CDF:
$F(\phi_k)=u_k$, where $F(\phi)=\int_0^\phi\rho(t)\dd t$.
The RoPE inverse frequencies are $\omega_k=\omega(\phi_k)=b^{-\phi_k}$.
\end{definition}

\begin{remark}[Discrete endpoint convention]
Under $u_k=k/N$, the discrete grid does not include $u=1$ (hence typically
$\phi_{N-1}<1$). We use the continuous interval $[0,1]$ for analysis and treat
$\phi=1$ as a normalized low-frequency boundary.
\end{remark}

\section{Exact Kernel and Broadband Projection}

\subsection{Exact kernel for the power-law prior}

\begin{definition}[Power-law distance prior]
For a training context length $L>1$, define
\begin{equation}
D(\Delta)=\frac{1}{\Delta\ln L}\,\mathbf{1}_{\Delta\in[1,L]}.
\end{equation}
\end{definition}

\begin{definition}[Phase-collision kernel]
The exact phase-collision kernel is the integral operator kernel
\begin{equation}
K(\phi_1,\phi_2)=\int_1^L D(\Delta)\cos(\omega(\phi_1)\Delta)\cos(\omega(\phi_2)\Delta)\dd\Delta.
\label{eq:exact_kernel_def}
\end{equation}
\end{definition}

Using $\cos a\cos b=\tfrac{1}{2}(\cos(a-b)+\cos(a+b))$ and the cosine integral
$\Ci(x)=-\int_x^\infty \cos t/t\dd t$, we obtain the closed form
\begin{equation}
K(\phi_1,\phi_2)
=\frac{1}{2\ln L}\Bigl[\Ci(\omega_- L)-\Ci(\omega_-)+\Ci(\omega_+ L)-\Ci(\omega_+)\Bigr],
\label{eq:exact_kernel_ci}
\end{equation}
where $\omega_-=\lvert \omega(\phi_1)-\omega(\phi_2)\rvert$ and
$\omega_+=\omega(\phi_1)+\omega(\phi_2)$.

\subsection{Broadband projection (the only approximation)}

\begin{definition}[Broadband projection surrogate]
We approximate the exact kernel by the two-parameter family
\begin{equation}
K_{\mathrm{app}}(\phi_1,\phi_2)=\alpha\,\delta(\phi_1-\phi_2)+\beta\,\min(\phi_1,\phi_2),
\label{eq:broadband_projection}
\end{equation}
where $\alpha,\beta>0$ are constants (fit/projection coefficients) and the
residual is $\E=K-K_{\mathrm{app}}$.
\end{definition}

\begin{remark}[What is and is not claimed]
Equation~\eqref{eq:broadband_projection} is the only approximation in the EVQ-Cosh
derivation chain. In this project, $(\alpha,\beta)$ are treated as broadband
effective coefficients (e.g. via Hilbert--Schmidt projection / finite-$N$
calibration), and we do \emph{not} claim the pointwise residual $\E$ is
asymptotically small for large architectural bases (e.g.\ $b\approx 5\times10^5$).
What is supported empirically is that the mid-band structure is captured with
high fidelity ($R^2>0.9$ under the empirical distance prior from FineWeb-Edu; $R^2\approx 0.96$ at $L=512$, declining to $\approx 0.87$ at $L=16\mathrm{K}$ as the RoPE collision bottleneck intensifies) and that the remaining matrix
residual is dominated by boundary and diagonal-ridge artifacts.
\end{remark}

\section{Operator Viewpoint: \texorpdfstring{$\min$}{min} as a Green's Function}

\begin{definition}[Mixed-boundary Laplacian]
Let $A=-\frac{\mathrm{d}^2}{\mathrm{d}\phi^2}$ on $[0,1]$ with mixed boundary
conditions $g(0)=0$ and $g'(1)=0$. Its inverse $A^{-1}$ is the integral operator
with kernel $\min(\phi_1,\phi_2)$.
\end{definition}

\begin{lemma}[Green's function identity]
For fixed $\psi\in(0,1)$, the function $G(\phi,\psi)=\min(\phi,\psi)$ satisfies
\begin{equation}
-\frac{\partial^2}{\partial \phi^2}G(\phi,\psi)=\delta(\phi-\psi),\qquad
G(0,\psi)=0,\qquad
\frac{\partial}{\partial\phi}G(1,\psi)=0,
\end{equation}
in the sense of distributions.
\end{lemma}
\begin{proof}
For $\phi<\psi$, $G(\phi,\psi)=\phi$; for $\phi>\psi$, $G(\phi,\psi)=\psi$.
Hence $\partial_\phi G$ jumps from $1$ to $0$ at $\phi=\psi$, so
$\partial_\phi^2 G=-\delta(\phi-\psi)$. The boundary conditions follow from the
piecewise form.
\end{proof}

\begin{remark}[Surrogate operator family]
With this convention, the broadband surrogate
$K_{\mathrm{app}}=\alpha I+\beta A^{-1}$ is the sum of an identity operator and a
resolvent (Green's) operator. This is a convenient continuous model for the
discrete diagonal ridge ($I$) plus smooth off-diagonal covariance ($A^{-1}$).
\end{remark}

\section{Variational Functional and Exact ODE Solution}

\subsection{Broadband functional}

Consider the broadband functional (continuous density formulation)
\begin{equation}
\begin{aligned}
\mathcal{J}[\rho]
&=\frac{\alpha}{2}\int_0^1\rho(\phi)^2\,\dd\phi
 +\frac{\beta}{2}\int_0^1\!\!\int_0^1\rho(\phi_1)\rho(\phi_2)\min(\phi_1,\phi_2)\,\dd\phi_1\dd\phi_2\\
&\quad-\mu_F\int_0^1\rho(\phi)b^{-2\phi}\,\dd\phi
 +\lambda\!\left(\int_0^1\rho(\phi)\,\dd\phi-1\right).
\end{aligned}
\label{eq:broadband_functional}
\end{equation}
Here $\mu_F$ controls the strength of the Fisher-resolution term and $\lambda$ enforces
normalization.

\begin{theorem}[Exact Euler--Lagrange ODE and general solution]
Let $\rho^\star$ be a stationary point of~\eqref{eq:broadband_functional}. Then
$\rho^\star$ satisfies the integral stationarity condition
\begin{equation}
\alpha\rho(\phi)+\beta\int_0^1\rho(\psi)\min(\phi,\psi)\,\dd\psi-\mu_F b^{-2\phi}+\lambda=0.
\label{eq:stationarity_integral}
\end{equation}
Moreover, $\rho$ satisfies the (interior) non-homogeneous ODE
\begin{equation}
\rho''(\phi)-\tau^2\rho(\phi)=\gamma_F b^{-2\phi},\qquad
\tau=\sqrt{\beta/\alpha},\qquad
\gamma_F=\frac{\mu_F(2\ln b)^2}{\alpha},
\label{eq:rho_ode}
\end{equation}
with boundary condition $\rho'(1)=-(2\mu_F\ln b/\alpha)b^{-2}$.
For $\tau\neq 2\ln b$, the general solution is
\begin{equation}
\rho^\star(\phi)=C_1\cosh(\tau\phi)+C_2\sinh(\tau\phi)+P b^{-2\phi},\qquad
P=\frac{\gamma_F}{4(\ln b)^2-\tau^2},
\label{eq:rho_general_solution}
\end{equation}
with constants $C_1,C_2$ determined by $\rho'(1)$ and $\int_0^1\rho=1$.
\end{theorem}
\begin{proof}
Taking the first variation of~\eqref{eq:broadband_functional} yields
Equation~\eqref{eq:stationarity_integral}. Differentiating once in $\phi$ and
using $\frac{\mathrm{d}}{\mathrm{d}\phi}\int_0^1\rho(\psi)\min(\phi,\psi)\dd\psi
=\int_\phi^1\rho(\psi)\dd\psi$ gives
\begin{equation}
\alpha\rho'(\phi)+\beta\int_{\phi}^{1}\rho(\psi)\,\dd\psi + 2\mu_F\ln b\;b^{-2\phi}=0.
\end{equation}
Evaluating at $\phi=1$ yields the boundary condition on $\rho'(1)$.
Differentiating again and using $\frac{\mathrm{d}}{\mathrm{d}\phi}\int_\phi^1\rho(\psi)\dd\psi=-\rho(\phi)$
gives $\alpha\rho''(\phi)-\beta\rho(\phi)=\mu_F(2\ln b)^2 b^{-2\phi}$, i.e.\
\eqref{eq:rho_ode}. The closed-form solution~\eqref{eq:rho_general_solution}
follows from constant-coefficient ODE theory; a particular solution $Pb^{-2\phi}$
exists whenever $\tau\neq 2\ln b$.
\end{proof}

\subsection{The cosh tether as the dominant broadband component}

\begin{remark}[Dominant cosh component]
In practice, $\tau=\mathcal{O}(1)$ while $2\ln b\approx 20$--$30$ for common
architectural bases, so the Fisher pulse $P b^{-2\phi}$ is sharply localized
near $\phi=0$. The macroscopic interference-suppression shape is therefore
captured by the homogeneous ``cosh tether''.
\end{remark}

\begin{corollary}[Pure tether model yields an exact cosh density]
Consider the simplified (Tikhonov + Brownian) functional with Fisher term removed:
\begin{equation}
\mathcal{J}_{\alpha,\beta}[\rho]=\frac{\alpha}{2}\int_0^1\rho(\phi)^2\,\dd\phi
\;+\;\frac{\beta}{2}\int_0^1\!\!\int_0^1\rho(\phi_1)\rho(\phi_2)\min(\phi_1,\phi_2)\,\dd\phi_1\dd\phi_2
\;+\;\lambda\!\left(\int_0^1\rho(\phi)\,\dd\phi-1\right).
\label{eq:J_alpha}
\end{equation}
Then the unique minimizer satisfies
\begin{equation}
\rho_{\alpha,\beta}^\star(\phi)=\frac{\tau\cosh(\tau(1-\phi))}{\sinh\tau},\qquad \tau=\sqrt{\beta/\alpha}.
\label{eq:rho_cosh_density}
\end{equation}
\end{corollary}
\begin{proof}
The stationarity equation is $\alpha\rho(\phi)+\beta\int_0^1\rho(\psi)\min(\phi,\psi)\dd\psi+\lambda=0$.
Differentiating twice as in the theorem proof yields $\alpha\rho''(\phi)-\beta\rho(\phi)=0$
with boundary condition $\rho'(1)=0$. Writing $\tau=\sqrt{\beta/\alpha}$, the general solution is
$\rho(\phi)=A\cosh(\tau(1-\phi))+B\sinh(\tau(1-\phi))$; the boundary condition forces $B=0$.
Normalization fixes $A=\tau/\sinh\tau$.
\end{proof}

\section{Exact Variational Quantization (EVQ): Closed-Form Warp}

\begin{theorem}[CDF and inverse CDF for the cosh density]
Let $\rho_\tau(\phi)=\frac{\tau\cosh(\tau(1-\phi))}{\sinh\tau}$ for $\tau>0$.
Then its CDF is
\begin{equation}
F_\tau(\phi)=\int_0^\phi\rho_\tau(t)\dd t
=1-\frac{\sinh(\tau(1-\phi))}{\sinh\tau}.
\label{eq:cosh_cdf}
\end{equation}
For $u\in(0,1)$, the inverse CDF is
\begin{equation}
F_\tau^{-1}(u)=\phi(u;\tau)=1-\frac{1}{\tau}\asinh\!\bigl((1-u)\sinh\tau\bigr).
\label{eq:evq_warp}
\end{equation}
\end{theorem}
\begin{proof}
Integrate $\cosh(\tau(1-\phi))$ explicitly:
$\int_0^\phi \cosh(\tau(1-t))\dd t
=\frac{1}{\tau}(\sinh\tau-\sinh(\tau(1-\phi)))$.
Multiplying by $\tau/\sinh\tau$ yields~\eqref{eq:cosh_cdf}. Inverting
$1-u=\sinh(\tau(1-\phi))/\sinh\tau$ gives~\eqref{eq:evq_warp}.
\end{proof}

\begin{corollary}[Geometric is the $\tau\to 0$ limit]
Let $u_k=k/N$. Then $\phi_k(\tau)=\phi(u_k;\tau)$ satisfies
$\lim_{\tau\to 0}\phi_k(\tau)=u_k$. Equivalently, EVQ degrades to geometric RoPE
as $\tau\to 0$.
\end{corollary}
\begin{proof}
For small $\tau$, $\sinh\tau=\tau+\mathcal{O}(\tau^3)$ and
$\asinh(x)=x+\mathcal{O}(x^3)$. Hence
\[
\phi_k(\tau)=1-\frac{1}{\tau}\asinh((1-u_k)\tau+\mathcal{O}(\tau^3))
=1-\frac{(1-u_k)\tau+\mathcal{O}(\tau^3)}{\tau}
=u_k+\mathcal{O}(\tau^2).
\]
\end{proof}

\begin{proposition}[Endpoint anchoring]
For any $\tau>0$, the EVQ warp satisfies $\phi(0;\tau)=0$ and $\phi(1;\tau)=1$.
In particular, under the discrete choice $u_0=0$ we have $\phi_0=0$ and
$\omega_0=b^{-\phi_0}=1$ exactly, matching the geometric highest-frequency
channel.
\end{proposition}
\begin{proof}
Immediate from~\eqref{eq:evq_warp}: $\asinh(\sinh\tau)=\tau$ and $\asinh(0)=0$.
\end{proof}

\begin{remark}[Local expansion (checked)]
Expanding~\eqref{eq:evq_warp} at $\tau=0$ gives
\begin{equation}
\phi(u;\tau)=u-\frac{u(1-u)(2-u)}{6}\tau^2+\mathcal{O}(\tau^4),
\label{eq:phi_taylor}
\end{equation}
so for $\tau>0$ we have $\phi(u;\tau)<u$ for $u\in(0,1)$ (a shift toward the
high-frequency end), while the \emph{spacing} is reallocated asymmetrically:
low-frequency gaps increase and high-frequency gaps shrink.
\end{remark}

\section{Waterbed Inequality (Integrated Trade-off)}

\begin{proposition}[Jensen waterbed bound]
Define local Fisher information $\mathcal{I}(\phi)=c\,\rho(\phi)\,b^{-2\phi}$ for
some constant $c>0$, and assume a local error proxy satisfies $E(\phi)\ge 1/\mathcal{I}(\phi)$.
Then
\begin{equation}
\int_0^1 \ln E(\phi)\,\dd\phi \ge \ln b-\ln c.
\label{eq:waterbed}
\end{equation}
Moreover, if $E(\phi)=1/\mathcal{I}(\phi)$ and $\rho$ is unconstrained beyond
normalization, equality holds if and only if $\rho$ is constant (geometric).
\end{proposition}
\begin{proof}
By $E\ge 1/\mathcal{I}$, $\ln E\ge -\ln \mathcal{I}$. Hence
\[
\int_0^1 \ln E\,\dd\phi \ge -\int_0^1 \ln c\,\dd\phi-\int_0^1\ln\rho(\phi)\dd\phi +2\ln b\int_0^1 \phi\,\dd\phi.
\]
Since $\int_0^1\phi\,\dd\phi=\tfrac{1}{2}$, this is
$\int_0^1\ln E\ge -\ln c-\int_0^1\ln\rho +\ln b$.
Finally, Jensen's inequality for the convex function $-\ln x$ gives
$-\int_0^1\ln\rho(\phi)\dd\phi \ge -\ln\!\int_0^1\rho(\phi)\dd\phi = 0$,
with equality iff $\rho$ is constant.
\end{proof}

\section{Scaling Law for \texorpdfstring{$\tau^\star$}{tau*} (Status: Conjecture)}

\begin{remark}[What is proved vs.\ what is fit]
The project contains strong empirical evidence that the optimal EVQ parameter
$\tau^\star$ scales approximately as $d_{\mathrm{head}}/\sqrt{L}$ at fixed base,
but a fully rigorous derivation of the $\alpha(L,b)$ calibration is still an
open proof task. We therefore record the scaling as a conjecture.
\end{remark}

\begin{proposition}[Single-variable scaling law (conjectural statement)]
Fix $b$ and $d_{\mathrm{head}}$. For a training length $L$, the optimal EVQ
parameter obeys
\begin{equation}
\tau^\star(L)\approx \frac{d_{\mathrm{head}}}{\sqrt{L}}.
\label{eq:tau_scaling_law}
\end{equation}
\end{proposition}

\section{Summary: Copy-Paste Formula Sheet}

\paragraph{EVQ warp (inverse CDF).}
For $N=d_{\mathrm{head}}/2$ and $u_k=k/N$:
\begin{equation}
\phi_k(\tau)=1-\frac{1}{\tau}\asinh\!\bigl((1-u_k)\sinh\tau\bigr),
\qquad
\omega_k=b^{-\phi_k}.
\end{equation}

\paragraph{Cosh density.}
\begin{equation}
\rho_\tau(\phi)=\frac{\tau\cosh(\tau(1-\phi))}{\sinh\tau},
\qquad
F_\tau(\phi)=1-\frac{\sinh(\tau(1-\phi))}{\sinh\tau}.
\end{equation}

\paragraph{Broadband ODE and general solution.}
\begin{equation}
\rho''-\tau^2\rho=\gamma_F b^{-2\phi},\qquad
\rho(\phi)=C_1\cosh(\tau\phi)+C_2\sinh(\tau\phi)+\frac{\gamma_F}{4(\ln b)^2-\tau^2}b^{-2\phi}.
\end{equation}

\paragraph{Waterbed inequality.}
\begin{equation}
\int_0^1 \ln E(\phi)\,\dd\phi \ge \ln b-\ln c.
\end{equation}

\end{document}
